\color{unidarkgreen}{\subsection{Контроль положения выключателей}}
\color{black}

\begin{enumerate}[label=\arabic{section}.\arabic{subsection}.\arabic*,  align=left, labelsep=4pt, leftmargin=0pt, itemindent=57pt]

\item
Контроль положения выключателей низкой и средней стороны используется для пуска МТЗ ВН, для
функции КПОН и в логике отключения СВ СН.  

\begin{figure}[H]
\centering
\includegraphics[width=0.88\textwidth,height=0.88\textheight,keepaspectratio]{\fbpath/081.3451 КПВ/img/KPV.pdf}
\captionof{figure}{Функциональная схема контроля положения выключателей}\label{pict:KPV}
\end{figure}

\item
Предусмотрена возможность приема сигналов включенного и отключенного положения выключателей
ввода сторон НН (НН1), НН2, СН1 (НН3) и СН2 (НН4). 

\item
Ввод в работу контроля положения В каждой стороны осуществляется независимо для каждой стороны
с помощью уставок \textbf{«Контроль В1 НН1(НН)»}, \textbf{«Контроль В2 НН1(НН)»}, \textbf{«Контроль В1 НН2(СН)»} и
\textbf{«Контроль В2 НН2(СН)»}.  

\item
Оперативный ввод/вывод контроля положения В производится с помощью БУ (блоков управления
функциями) \textit{«Вывод В1 НН1(НН)»}, \textit{«Вывод В2 НН1(НН)»}, \textit{«Вывод В1 НН2(СН)»}, \textit{«Вывод В2 НН2 (СН)»}.


\small
\setlength{\extrarowheight}{0.05cm} % высота ячеек в таблице
\begin{longtable}
{|>{\centering\arraybackslash}m{4.5cm}|>{\centering\arraybackslash}m{4.5cm}|>{\centering\arraybackslash}m{3cm}|>{\centering\arraybackslash}m{3.5cm}|}
\caption{Перечень уставочных параметров контроля положения выключателей\hfill\vspace{-0.5\baselineskip}}\label{kpv:tbl1} \\ 
\hhline{|-|-|-|-|}
\rowcolor{gray!30}
Наименование уставки & Диапазон & Шаг &  Значение по умолчанию \\ 
\hhline{|-|-|-|-|}
\endfirsthead

\caption*{\hspace{3pt}\emph{Продолжение таблицы \ref{kpv:tbl1}\hfill\vspace{-0.5\baselineskip}}} \\ % сделано по ГОСТ 2.105 п.6.8.7
\hhline{|-|-|-|-|}
\rowcolor{gray!30}
Наименование уставки & Диапазон & Шаг &  Значение по умолчанию \\  
\hhline{|-|-|-|-|}
\endhead
\hhline{|-|-|-|-|}

\endfoot
\endlastfoot

%%%%%%%%%%%%содержимое%%%%%%%%%%%%%%

\end{longtable}
\setlength{\extrarowheight}{0.0cm} % возврат параметра высота ячеек в таблице
\normalsize










































\end{enumerate} 